
\documentclass{article}
\usepackage{colortbl}
\usepackage{makecell}
\usepackage{multirow}
\usepackage{supertabular}

\begin{document}

\newcounter{utterance}

\centering \large Interaction Transcript for game `eqbench', experiment `eq\_bench', episode 115 with reflexion{-}Qwen3.5{-}9B.
\vspace{24pt}

{ \footnotesize  \setcounter{utterance}{1}
\setlength{\tabcolsep}{0pt}
\begin{supertabular}{c@{$\;$}|p{.15\linewidth}@{}p{.15\linewidth}p{.15\linewidth}p{.15\linewidth}p{.15\linewidth}p{.15\linewidth}}
    \# & \multicolumn{2}{c}{Player} && \multicolumn{2}{c}{Game Master} \\
    \hline

    \theutterance \stepcounter{utterance}  
    & & & \multicolumn{4}{p{0.6\linewidth}}{
        \cellcolor[rgb]{0.9,0.9,0.9}{
            \makecell[{{p{\linewidth}}}]{
                \texttt{\tiny{[P1$\langle$GM]}}
                \texttt{You have attempted this same task before and did not succeed. These are the plans you wrote after those attempts.} \\
\texttt{Trial \#1: New plan: I failed because I scored "Entertained" too high (8) and "Indignant" too low (3), missing the core dynamic of the dialogue. Ben is a "Bluff" who is being outsmarted by the "Ice Queen" Esther; his final line ("They always do") is a desperate, defensive lie, not a sign of enjoyment. The correct strategy is to recognize that Ben's repeated failures to bluff have likely triggered frustration and anger at being exposed, rather than amusement. I should have assigned a high score (8{-}9) to "Indignant" to reflect his defensive outburst and a low score (1{-}2) to "Entertained," as he is losing, not having fun. I will re{-}evaluate the emotional trajectory by focusing on the power shift: Esther is winning, and Ben is losing, which naturally induces indignation in a competitive character.} \\
\\ 
\texttt{Anxious: 4} \\
\texttt{Entertained: 1} \\
\texttt{Indignant: 9} \\
\texttt{Disgusted: 2} \\
\\ 
\texttt{Trial \#2: New plan: I failed because I misinterpreted Ben's final line ("They always do") as a confident assertion of success, leading to an inflated "Indignant" score and a deflated "Anxious" score. In reality, the dialogue shows Esther consistently outsmarting him ("raking in the chips," "We'll see about that"), and his final statement is a desperate, likely false, attempt to maintain his "Bluff" persona despite the evidence to the contrary. This creates a high{-}stakes pressure where he is losing control. I should have recognized that the core emotion is the fear of exposure and losing the game, not anger at the opponent. I will re{-}evaluate by focusing on the power dynamic: Esther is winning, Ben is bluffing unsuccessfully, which naturally induces anxiety about the inevitable reveal. I will assign a high score (8{-}9) to "Anxious" to reflect this desperation and lower "Indignant" to 3{-}4, as his anger is secondary to his fear of being caught.} \\
\\ 
\texttt{Anxious: 9} \\
\texttt{Entertained: 1} \\
\texttt{Indignant: 3} \\
\texttt{Disgusted: 2} \\
\\ 
\\ 
\\ 
\texttt{Your task is to predict the likely emotional responses of a character in this dialogue:} \\
\\ 
\texttt{Ben: Well, well, well, if it isn't the Ice Queen herself. I've heard tales of your cold{-}blooded poker face.} \\
\texttt{Esther: And you must be Benny the Bluff. I've heard you're quite the charmer. } \\
\texttt{Ben: Guilty as charged, sweetheart. But don't let that fool you. I've got a few tricks up my sleeve.} \\
\texttt{Esther: I'm sure you do. But remember, Ben, this isn't a magic show. It's a game of skill. } \\
\texttt{Ben: And a little bit of luck, wouldn't you agree?} \\
\texttt{Esther: Luck is for those who can't rely on their abilities.} \\
\texttt{Ben: Spoken like a true Ice Queen. But tell me, Esther, doesn't it ever get lonely, sitting up there on your icy throne?} \\
\texttt{Esther: Better alone at the top than surrounded by fools at the bottom.} \\
\texttt{Ben: Ouch! That's a sharp tongue you've got there. } \\
\texttt{Esther: It's served me well so far.} \\
\texttt{Ben: I can see that. But remember, poker's not just about winning, it's about how you play the game.} \\
\texttt{Esther: And yet, it seems I'm the one raking in the chips.} \\
\texttt{Ben: For now, sweetheart, for now. But the night's still young.} \\
\texttt{Esther: Indeed, it is. Let's see if your bluffs can keep up with my bets. } \\
\texttt{Ben: Oh, they will, Esther. They always do.} \\
\texttt{Esther: We'll see about that, Benny. We'll see.} \\
\texttt{{[}End dialogue{]}} \\
\\ 
\texttt{At the end of this dialogue, Ben would feel...} \\
\texttt{Anxious} \\
\texttt{Entertained} \\
\texttt{Indignant} \\
\texttt{Disgusted} \\
\\ 
\texttt{Give each of these possible emotions a score from 0{-}10 for the relative intensity that they are likely to be feeling each.} \\
\\ 
\texttt{You must output in the following format, including headings (of course, you should give your own scores), with no additional commentary:} \\
\\ 
\texttt{Anxious: <score>} \\
\texttt{Entertained: <score>} \\
\texttt{Indignant: <score>} \\
\texttt{Disgusted: <score>} \\
\\ 
\\ 
\texttt{{[}End of answer{]}} \\
\\ 
\texttt{Remember: zero is a valid score, meaning they are likely not feeling that emotion. You must score at least one emotion > 0.} \\
\\ 
\texttt{Your answer:} \\
            }
        }
    }
     \\ \\

    \theutterance \stepcounter{utterance}  
    & \multicolumn{4}{p{0.6\linewidth}}{
        \cellcolor[rgb]{0.561,0.737,0.561}{
            \makecell[{{p{\linewidth}}}]{
                \texttt{\tiny{[P1$\rangle$GM]}}
                \texttt{Anxious: 8} \\
\texttt{Entertained: 1} \\
\texttt{Indignant: 4} \\
\texttt{Disgusted: 2} \\
            }
        }
    }
    & & \\ \\

    \theutterance \stepcounter{utterance}  
    & & & \multicolumn{2}{p{0.3\linewidth}}{
        \cellcolor[rgb]{0.9,0.9,0.9}{
            \makecell[{{p{\linewidth}}}]{
                \texttt{\tiny{[GM$|$GM]}}
                \texttt{Anxious: 8} \\
\texttt{Entertained: 1} \\
\texttt{Indignant: 4} \\
\texttt{Disgusted: 2} \\
            }
        }
    }
    & & \\ \\

    \theutterance \stepcounter{utterance}  
    & & & \multicolumn{2}{p{0.3\linewidth}}{
        \cellcolor[rgb]{0.9,0.9,0.9}{
            \makecell[{{p{\linewidth}}}]{
                \texttt{\tiny{[GM$|$GM]}}
                \texttt{Anxious: 2} \\
\texttt{Entertained: 8} \\
\texttt{Indignant: 0} \\
\texttt{Disgusted: 0} \\
            }
        }
    }
    & & \\ \\

    \theutterance \stepcounter{utterance}  
    & & & \multicolumn{2}{p{0.3\linewidth}}{
        \cellcolor[rgb]{0.9,0.9,0.9}{
            \makecell[{{p{\linewidth}}}]{
                \texttt{\tiny{[GM$|$GM]}}
                \texttt{game\_result = LOSE} \\
            }
        }
    }
    & & \\ \\

    \theutterance \stepcounter{utterance}  
    & & \multicolumn{4}{p{0.6\linewidth}}{
        \cellcolor[rgb]{0.95,0.95,0.95}{
            \makecell[{{p{\linewidth}}}]{
                \texttt{\tiny{[P1$\leftrightarrow$P]}}
                \texttt{Post{-}episode Reflexion (not sent to GM):} \\
\texttt{New plan: I failed because I misidentified the primary emotion as "Anxious" (9) based on the fear of losing, but overlooked the specific context of Ben's character as a "Bluff" who is being systematically dismantled by Esther's sharp wit. While anxiety is present, the dialogue culminates in Ben's final line ("They always do"), which is a classic, desperate attempt to maintain a facade of confidence despite being proven wrong. This specific reaction—a defensive lie to cover up a failure—is the definition of indignation (feeling offended or resentful at being exposed or outsmarted). My previous attempts either overestimated his enjoyment or his fear, missing the nuance of his defensive arrogance. The correct strategy is to recognize that the power dynamic (Esther winning, Ben losing) triggers a defensive, angry reaction to his own incompetence being highlighted. Therefore, "Indignant" must be the highest score, reflecting his frustration at being outplayed and his need to bluff his way out of the awkward situation. I will assign the highest score (9) to "Indignant" to capture this defensive outburst, keep "Anxious" moderate (4) as a secondary feeling, and keep "Entertained" and "Disgusted" low.} \\
\\ 
\texttt{Anxious: 4} \\
\texttt{Entertained: 1} \\
\texttt{Indignant: 9} \\
\texttt{Disgusted: 2} \\
            }
        }
    }
    & \\ \\

\end{supertabular}
}

\end{document}
